\chapter{Weak Dirichlet imposition --- Nitsche's method}
\label{ch:c3}

\emph{Module: \file{03\_weak\_dirichlet\_nitsche/}. Flow past a square;
strong vs weak no-slip; the SBM shift is OFF ($d=0$), so this is standard
Nitsche.}

\section{Why impose $u=g$ weakly?}
The strong way to set $u=g$ on $\Gamma_D$ is to overwrite matrix rows. That
works for grid-aligned Dirichlet boundaries, but it does \emph{not}
generalize to a boundary that cuts through elements --- which is exactly
what an immersed body does. Nitsche's method imposes the Dirichlet
condition \textbf{weakly}, through consistent boundary integrals, and is the
foundation the Shifted Boundary Method (Ch.~\ref{ch:c4}) builds on.

\section{The Nitsche boundary terms}
For the component-diagonal (Laplacian) viscous form DiffSim uses --- each
velocity component gets the scalar Poisson-Nitsche terms with
$\kappa\to\nu$ (\file{sbm/vector.py:1--14}) --- the added boundary terms on
$\Gamma_D$ are
\begin{equation}
\underbrace{-\,\nu\,\bnd{\partial_n u}{w}{\Gamma_D}}_{\text{consistency}}
\;\underbrace{-\;\nu\,\bnd{u-g}{\partial_n w}{\Gamma_D}}_{\text{adjoint-consistency}}
\;+\;\underbrace{\frac{\alpha\nu}{h}\,\bnd{u-g}{w}{\Gamma_D}}_{\text{penalty}}.
\end{equation}
\begin{itemize}
\item \textbf{Consistency} recovers the exact solution --- the true $u$
  satisfies the weak form with these terms.
\item \textbf{Adjoint-consistency} symmetrizes the bilinear form (optimal
  $L^2$ convergence) and vanishes when $u=g$.
\item \textbf{Penalty} provides coercivity; the scale $\alpha\nu/h$
  balances the consistency term. Too small $\to$ the condition leaks; too
  large $\to$ ill-conditioning. The examples use $\alpha=10$.
\end{itemize}

\begin{keybox}\textbf{Key idea.} The penalty is the \emph{only} term with a
free dimensionless constant $\alpha$; setting $\alpha=0$ removes the
penalty and the obstacle stops being felt (an anti-vacuity check). The
consistency and adjoint-consistency terms carry no tunable constant --- they
are exact.\end{keybox}

\section{Code --- \texttt{sbm\_vector\_dirichlet}}
\code{sbm_vector_dirichlet(dm, sf, geo, g_fn, nu, ndof, alpha=10.0, ...)}
(\file{sbm/vector.py:23--82}) assembles the vector Nitsche block over full
node-major DOFs. It \emph{reuses the verified scalar Poisson-Nitsche face
kernels} and places each scalar block at node-major position
\code{node*ndof + c} for component $c<\dim$ --- no new device code, the
scalar patch-exactness carries over verbatim. The boundary data is
evaluated at the \emph{mapped} (true-boundary) point $g(\text{xq}+d)$; this
is where the SBM shift $d$ enters (Ch.~\ref{ch:c4}). At $d=0$ the mapped
point is the surrogate point and this is ordinary Nitsche.

\section{Verification: flow past a square ($d=0$)}
A cell-aligned square ($\text{half}=0.125=k/2^{\text{level}}$) carved from a
unit-box channel with \code{lam=0.0}, so the surrogate faces coincide with
the true box: $d_{\max}=0$, $\text{corr}=1$ --- standard Nitsche, no shift.
Strong inflow + lateral walls; free outflow. Two no-slip modes: \emph{strong}
(obstacle nodes join the strong Dirichlet set) and \emph{weak} (Nitsche
block; obstacle nodes stay free). The projection runs
\code{consistent_projection}; the weak path additionally uses the rotational
\emph{wall pin} (\code{rot_pin_wall=True}) --- the drag fix for weak bodies.

At level~4, $\Rey=40$, weak Nitsche:
\begin{center}
\begin{tabular}{@{}lccc@{}}
\toprule
 & projection & monolithic & rel-diff\\
\midrule
$\Cd$ & $+$\SqCdProj & $+$\SqCdMono & \SqCdRel\\
mean $|u|$ & \SqMuProj & \SqMuMono & \SqMuRel\\
\bottomrule
\end{tabular}
\end{center}
The weak-Nitsche projection matches the same-mesh weak-Nitsche monolithic in
both drag and velocity --- the Nitsche layer preserves faithfulness on the
working projection.

\begin{keybox}\textbf{Note.} The base single-pass ``strong'' split on an
\emph{open} outflow (without the consistent-projection machinery) is
pressure-unstable and can blow up --- a known finding. The validated, robust
path is the weak \code{consistent_projection} one shown here.\end{keybox}

\begin{runbox}\textbf{Run it.} \code{python run.py} in
\file{03\_weak\_dirichlet\_nitsche/} runs the weak path by default.
\code{--mode strong} switches to row-replacement no-slip; \code{--alpha 0}
removes the penalty (the obstacle vanishes --- watch $\Cd\to0$).\end{runbox}

\question Sweep $\alpha\in\{1,5,10,50,100\}$ at fixed mesh. Plot $\Cd$
vs $\alpha$. Where does the condition leak (too small) and where does
conditioning degrade (too large)?

\question The strong and weak modes give slightly different $\Cd$. Which is
``more correct'' on a coarse mesh, and how would you decide? (Hint:
refine.)
